\chapter{簇与理想}\label{chap:varieties-and-ideals}

\noindent\textbf{第 \thechapter{} 章学习目标}
\begin{itemize}
  \item 定义仿射空间 $\mathbb{A}^{n}_{K}$, 并能在其中进行运算.
  \item 使用 $\mathcal{V}$ 和 $\mathcal{I}$ 算子, 在 $K[x_{1},\ldots,x_{n}]$ 的子集与 $\mathbb{A}^{n}_{K}$ 的子集之间转换.
  \item 定义 $\mathbb{A}^{n}_{K}$ 中的仿射簇和 $K[x_{1},\ldots,x_{n}]$ 中的根理想, 并举例说明.
  \item 叙述零点定理, 并用它描述仿射簇与根理想之间的双射.
\end{itemize}

从根本上说, 代数几何是一本在两种语言之间进行翻译的词典: 一种是代数的语言, 另一种是几何的语言. 和任何双语词典一样, 它可以在两个方向上进行翻译. 给定一个代数对象, 例如一个多项式, 我们通过确定它的零点集得到一个几何对象. 例如, 多项式 $x^{2}+y^{2}-1$ 的零点集是平面上的单位圆. 反过来, 给定一个几何对象, 我们通过确定在给定几何集合上处处为零的所有多项式所成的集合, 得到一个代数对象.

在这一章, 我们借助 $\mathcal{V}$ 和 $\mathcal{I}$ 算子精确表述这两种运算, 正式开始学习代数几何. 关键在于, 我们发现这些算子不是满射. 并非每个几何集合都是某族多项式的零点集; 能以这种方式得到的几何集合称为仿射簇. 反过来, 并非每个多项式集合都能从某个几何集合出发, 通过求出在其上处处为零的多项式得到. 研究能够以这种方式得到的代数集合, 将引出根理想这一代数概念.

本章最后将叙述一个结果. 它本可以恰当地称为\textquotedbl{}代数几何基本定理\textquotedbl{}, 不过传统上它用德语名称 Nullstellensatz (零点定理). 在一个关键假设下, 即基域是代数闭域, 零点定理断言: 若几何一侧限于仿射簇, 代数一侧限于根理想, 那么 $\mathcal{V}$ 和 $\mathcal{I}$ 算子就在代数与几何之间给出了一本真正的词典, 用数学的话说, 给出了一个双射.

\par\smallskip
{\footnotesize
\noindent\textcopyright{} The Author(s) 2025\par
\noindent E. Clader, D. Ross, \emph{Beginning in Algebraic Geometry}, Undergraduate Texts in Mathematics,
\href{https://doi.org/10.1007/978-3-031-88819-9_2}{\nolinkurl{https://doi.org/10.1007/978-3-031-88819-9_2}}.\par}
